% !TEX root = task-1-14.tex
% Editable task bank. Append new BankTask blocks at the end of the document.
\documentclass[12pt,a4paper]{article}
\usepackage{physics-textbook}
\usepackage{graphicx,needspace}
\usetikzlibrary{decorations.pathmorphing}
\newif\ifShowSolutions
\ifdefined\TutorEdition
  \ShowSolutionstrue
\else
  \ShowSolutionsfalse % Change to \ShowSolutionstrue to display solutions.
\fi
\newcounter{banktask}
\newsavebox{\BankTaskBox}
\newcommand{\BankTaskContent}[3]{%
  \noindent\textbf{\thebanktask.}\quad #2\par
  \ifShowSolutions
    \Needspace{50mm}% Keep the solution heading with its opening text or diagram.
    #3\par
    \ifdefined\TutorEdition
      \ifcsname TutorSource#1\endcsname
        \nopagebreak\smallskip{\small\csname TutorSource#1\endcsname\par}%
      \fi
    \fi
  \fi
}
\newcommand{\BankTask}[3]{%
  \refstepcounter{banktask}\par\medskip
  \sbox{\BankTaskBox}{\begin{minipage}{\linewidth}%
    \BankTaskContent{#1}{#2}{#3}\end{minipage}}%
  \ifdim\dimexpr\ht\BankTaskBox+\dp\BankTaskBox\relax>\dimexpr\textheight-20mm\relax
    \Needspace{35mm}\BankTaskContent{#1}{#2}{#3}%
  \else
    \Needspace{\dimexpr\ht\BankTaskBox+\dp\BankTaskBox\relax}%
    \noindent\usebox{\BankTaskBox}\par
  \fi
  \medskip
}
\begin{document}
\BookSubsection{1.14}{Импульс}


% Source F579FA: https://ege.fipi.ru/bank/questions.php?search=1&proj=BA1F39653304A5B041B656915DC36B38&qid=F579FA
\BankTask{F579FA}{
Отношение скорости грузовика к скорости легкового автомобиля $\frac{\upsilon_1}{\upsilon_2}=\frac{1}{2}.$ Масса грузовика $m_1=3000\,\text{кг}$. Какова масса легкового автомобиля, если отношение импульса грузовика к импульсу легкового автомобиля равно $1{,}5$? Результат выразите в $\text{кг}$.
}{
\smallskip\noindent\textbf{Решение от GPT}\par
$p_1/p_2=(m_1/m_2)(v_1/v_2)$. Поэтому $m_2=m_1(v_1/v_2)/(p_1/p_2)=3000\cdot0{,}5/1{,}5=1000\,\text{кг}$.
}

% Source 72C7F7: https://ege.fipi.ru/bank/questions.php?search=1&proj=BA1F39653304A5B041B656915DC36B38&qid=72C7F7
\BankTask{72C7F7}{
Отношение импульса легкового автомобиля к импульсу мотоцикла $\frac{p_1}{p_2}=5$. Каково отношение их скоростей $\frac{\upsilon_1}{\upsilon_2},$ если отношение массы легкового автомобиля к массе мотоцикла $\frac{m_1}{m_2}=2{,}5$?
}{
\smallskip\noindent\textbf{Решение от GPT}\par
$v_1/v_2=(p_1/p_2)/(m_1/m_2)=5/2{,}5=2$.
}

% Source D82e7B: https://ege.fipi.ru/bank/questions.php?search=1&proj=BA1F39653304A5B041B656915DC36B38&qid=D82e7B
\BankTask{D82e7B}{
Два тела, каждое массой $2\,\text{кг}$, движутся на одной прямой, вдоль которой направлена ось $O x$ На рисунке приведены графики зависимости проекций $\upsilon_x$ их скорости от времени $t.$
\begin{center}\begin{tikzpicture}[x=.225cm,y=.45cm,>=stealth]
\draw[thin,dashed,xstep=2,ystep=1] (0,0) grid (20,8);
\draw[->] (0,0)--(22,0) node[right] {$t,\,\text{с}$};
\draw[->] (0,0)--(0,9) node[above] {$\upsilon_x,\,\text{м}/\text{с}$};
\foreach \x in {0,4,8,12,16} \node[below] at (\x,0) {$\x$};
\foreach \y in {2,4,6} \node[left] at (0,\y) {$\y$};
\draw[very thick] (0,7)--(20,{19/7}); \node at (6,6.8) {$1$};
\draw[very thick] (0,.5)--(20,5.5); \node at (8,1.7) {$2$};
\end{tikzpicture}\end{center}
\begin{enumerate}
\item Равнодействующая сил, действующих на тело $1$, остаётся постоянной в течение всего времени наблюдения.
\item В промежутке времени от $0$ до $14\,\text{с}$ тела двигались навстречу друг другу.
\item Путь, пройденный телом $1$ за промежуток времени от $0$ до $18\,\text{с}$, больше пути, пройденного телом $2$ за тот же промежуток времени.
\item Кинетическая энергия тела $2$ за промежуток времени от $6$ до $14\,\text{с}$ увеличилась в $2$ раза.
\item Модуль импульса тела $2$ в момент времени $t=10\,\text{с}$ равен $10\,\text{кг}\cdot\text{м}/\text{с}$.
\end{enumerate}
}{
\smallskip\noindent\textbf{Решение от GPT}\par
Линейные графики задают $v_{1x}=7-\frac{3}{14}t$ и $v_{2x}=0{,}5+0{,}25t$ в СИ. Равнодействующая на тело $1$ постоянна; обе скорости положительны. За первые $18\,\text{с}$ пути $s_1=7\cdot18-\frac{3}{28}\cdot18^2=\frac{639}{7}\,\text{м}\approx91{,}3\,\text{м}$ и $s_2=(0{,}5+5)18/2=49{,}5\,\text{м}$, значит $s_1>s_2$. Для тела $2$ $E_{\text{к}}(14)/E_{\text{к}}(6)=(4/2)^2=4$, $p(10)=2\cdot3=6\,\text{кг}\cdot\text{м}/\text{с}$. Ответ: $13$.
}

% Source D4E9EE: https://ege.fipi.ru/bank/questions.php?search=1&proj=BA1F39653304A5B041B656915DC36B38&qid=D4E9EE
\BankTask{D4E9EE}{
В инерциальной системе отсчёта вдоль оси $Ox$ движется тело массой $20\,\text{кг}$. На рисунке приведён график зависимости проекции скорости $V_{x}$ этого тела от времени $t$.
\begin{center}
\begin{tikzpicture}[x=0.06cm,y=0.6cm,>=stealth]
\draw[thin,black!65,xstep=10,ystep=1] (0,0) grid (100,4);
\draw[->] (0,0)--(110,0) node[right] {$t,\,\text{с}$};
\draw[->] (0,0)--(0,5) node[above] {$v_x,\,\text{м}/\text{с}$};

\node[below] at (10,0) {$10$};
\node[below] at (20,0) {$20$};
\node[below] at (30,0) {$30$};
\node[below] at (40,0) {$40$};
\node[below] at (50,0) {$50$};
\node[below] at (60,0) {$60$};
\node[below] at (70,0) {$70$};
\node[below] at (80,0) {$80$};
\node[below] at (90,0) {$90$};
\node[below] at (100,0) {$100$};
\node[left] at (0,0) {$0$};
\node[left] at (0,1) {$1$};
\node[left] at (0,2) {$2$};
\node[left] at (0,3) {$3$};
\node[left] at (0,4) {$4$};
\draw[very thick] (0,0)--(20,4)--(60,4)--(80,3)--(100,0);
\end{tikzpicture}
\end{center}
\begin{enumerate}
\item Модуль ускорения тела в промежутке времени от $0$ до $20\,\text{с}$ в два раза больше модуля ускорения тела в промежутке времени от $60$ до $80\,\text{с}$
\item В промежутке времени от $0$ до $10\,\text{с}$ тело переместилось на $20\,\text{м}$
\item В момент времени $40\,\text{с}$ равнодействующая сил, действующих на тело, равна $0$
\item В промежутке времени от $80$ до $100\,\text{с}$ импульс тела уменьшился на $60\,\text{кг}\cdot\text{м}/\text{с}$
\item Кинетическая энергия тела в промежутке времени от $10$ до $20\,\text{с}$ увеличилась в $2$ раза
\end{enumerate}
}{
\smallskip\noindent\textbf{Решение от GPT}\par
$|a_{0-20}|=4/20=0{,}2\,\text{м}/\text{с}^2$, $|a_{60-80}|=(4-3)/20=0{,}05\,\text{м}/\text{с}^2$ (отношение $4$). За первые $10\,\text{с}$ путь $2\cdot10/2=10\,\text{м}$. При $t=40\,\text{с}$ $a=0$ и $F=0$. За $80\! -\!100\,\text{с}$ импульс уменьшился на $20(3-0)=60\,\text{кг}\cdot\text{м}/\text{с}$. При $10\! -\!20\,\text{с}$ скорость удваивается, энергия увеличивается вчетверо. Ответ: $34$.
}

% Source DB3508: https://ege.fipi.ru/bank/questions.php?search=1&proj=BA1F39653304A5B041B656915DC36B38&qid=DB3508
\BankTask{DB3508}{
Шайба массой $m$ съезжает без трения из состояния покоя с вершины горки. У подножия горки потенциальная энергия шайбы равна нулю, а модуль её импульса равен $p$. Каковы высота горки и потенциальная энергия шайбы на её вершине? Ускорение свободного падения равно $g$.
\TaskMatchingColumns[.53\linewidth]{Физические величины}{
А) высота горки\par
Б) потенциальная энергия шайбы на вершине горки
}{Формулы}{
1) $\frac{p^2}{2m^2g}$\par
2) $\frac{p^2}{2mg}$\par
3) $\frac{p^2}{2m}$\par
4) $\frac{mp^2}{2g}$\par
}\begin{center}\begin{tabular}{|c|c|}\hline А & Б \\\hline\rule{0pt}{1.2em}\hspace{1cm} & \hspace{1cm} \\\hline\end{tabular}\end{center}
}{
\smallskip\noindent\textbf{Решение от GPT}\par
$mgh=mv^2/2=p^2/(2m)$. Поэтому $h=p^2/(2m^2g)$ (формула $1$), $E_{\text{п}}=p^2/(2m)$ (формула $3$). Ответ: $13$.
}

% Source 589e07: https://ege.fipi.ru/bank/questions.php?search=1&proj=BA1F39653304A5B041B656915DC36B38&qid=589e07
\BankTask{589e07}{
Мяч подбросили вверх с крыши сарая. Как изменяются по мере подъёма модуль импульса мяча и его потенциальная энергия относительно поверхности Земли? Сопротивление воздуха не учитывать.
\begin{enumerate}
\item увеличивается
\item уменьшается
\item не изменяется
\end{enumerate}
\TaskChangeTable{Модуль импульса мяча}{Потенциальная энергия мяча}
}{
\smallskip\noindent\textbf{Решение от GPT}\par
При подъёме скорость уменьшается, поэтому $p=mv$ уменьшается. Высота растёт, значит $E_{\text{п}}=mgh$ увеличивается. Ответ: $21$.
}

% Source 9A9FD1: https://ege.fipi.ru/bank/questions.php?search=1&proj=BA1F39653304A5B041B656915DC36B38&qid=9A9FD1
\BankTask{9A9FD1}{
Спортсмен спускается на парашюте с постоянной скоростью. Как изменяются с течением времени в процессе спуска импульс спортсмена и его потенциальная энергия?
\begin{enumerate}
\item увеличивается
\item уменьшается
\item не изменяется
\end{enumerate}
\TaskChangeTable{Импульс}{Потенциальная энергия}
}{
\smallskip\noindent\textbf{Решение от GPT}\par
При постоянной скорости $\vec p=m\vec v$ не меняется. Высота уменьшается, поэтому $E_{\text{п}}=mgh$ уменьшается. Ответ: $32$.
}

\end{document}
